% Part: proof-theory
% Chapter: proof-search
% Section: completeness

\documentclass[../../../include/open-logic-section]{subfiles}

\begin{document}

\olfileid{pt}{ps}{cpl}

\olsection{Kelengkapan}

Saiki kita nuduhaké yèn algoritma panelusuran bukti bisa
dipercaya: yèn algoritma mandheg tanpa asil utawa ora tau
mandheg, ora ana !!{proof} tumrap sekuèn wiwitan
$\Gamma \Sequent \Delta$. Kanggo kuwi, kita netepaké
!!a{structure}~$\Struct M$ saéngga kanggo saben $!A \in \Gamma$,
$\Sat{M}{!A}$, lan kanggo saben~$!A \in \Delta$,
$\Sat/{M}{!A}$. Tegesé, kabèh !!{formula}s ing~$\Gamma$
bener lan kabèh !!{formula}s ing~$\Delta$ salah
ing~$\Struct M$; mula $\Gamma \Sequent \Delta$ ora valid.
Amarga \Log{G3c} sahih, $\Gamma \Sequent \Delta$ ora nduwèni
!!{proof}.

Kita netepaké $\Struct{M}$ saka pasangan $\Theta, \Xi$
kang isiné !!{formula}s lan himpunan~$C$ kang isiné
!!{constant}s. $\Theta, \Xi$, lan~$C$ dijupuk saka
\emph{cabang gagal} ing panelusuran bukti kang gagal
(mandheg tanpa asil utawa ora tau mandheg). Cabang gagal
iku rerangkéyan sekuèn $\Pi_n \Sequent \Lambda_n$ lan
himpunan !!{constant}s~$C_n$, kanthi
$\Pi_0 \Sequent \Lambda_0$ minangka sekuèn pungkasan
$\Gamma \Sequent \Delta$, lan
$\Pi_{n+1} \Sequent \Lambda_{n+1}$ minangka premis saka
$\Pi_n \Sequent \Lambda_n$ kang diasilaké ing tataran~$n+1$
nanging ora ditemokaké !!a{proof} déning algoritma.
Kanggo saben~$n$ kang nduwèni tataran salajengé, premis
kaya mangkono mesthi ana; yèn ora, algoritma wis
nemokaké !!a{proof}. Kita anggep $C_n$ minangka himpunan
!!{constant}s kang gegandhèngan karo~$\Pi_n
\Sequent \Lambda_n$.

Yèn algoritma mandheg tanpa asil ing sekuèn paling dhuwur
kang mung ngemot !!{formula}s atomik, cabang kang pungkasané
ing sekuèn iku dadi cabang gagal. Yèn algoritma ora tau
mandheg, algoritma terus ngasilaké wit-witan kang saya
gedhé. Wit-witan iku bisa dianggep minangka tahap-tahap
pambangunan wit tanpa wates (wit kang kawangun sawisé
algoritma lumaku tanpa wates). Wit iki nduwèni cabang
winates ing saben simpul: saben simpul mung nduwèni siji
utawa loro anak, yaiku sekuèn kang langsung ana ing
sadhuwuré. Miturut Léma K\H{o}nig, saben wit tanpa wates
kang cabang saben simpulé winates mesthi nduwèni cabang
tanpa wates. Cabang iki dadi cabang gagal yèn algoritma
panelusuran lumaku salawasé.

Yèn rerangkéyan $\Pi_n \Sequent \Lambda_n$ lan $C_n$
iku cabang gagal, tetepaké $\Theta = \bigcup_n \Pi_n$
lan $\Xi = \bigcup_n \Lambda_n$ (kanthi mbusak indeks lan
kedadéan ganda !!{formula}s), sarta $C = \bigcup_n C_n$.

Saiki kita bisa netepaké~$\Struct{M}$.
\begin{enumerate}
\item Domain $\Domain{M}$ yaiku himpunan tèrma kang bisa
dibangun saka~$C$ lan !!{function}s ing~$\Gamma \Sequent \Delta$
(yèn ana), nanging tanpa !!{variable}s.
\item Yèn $c \in \Domain{M}$, $\Assign{c}{M} = c$.
\item Yèn $f$ iku !!{function} aritas~$m$ ing
$\Gamma \Sequent \Delta$, lan $t_1$, \dots,
$t_m \in \Domain{M}$, mula
$\Assign{f}{M}(t_1, \dots, t_m) = f(t_1, \dots, t_m)$.
\item Yèn $R$ iku !!{predicate} aritas~$m$ ing
$\Gamma \Sequent \Delta$, mula
$\Assign{R}{M} = \Setabs{\tuple{t_1, \dots, t_m} \in
\Domain{M}^m}{R(t_1, \dots, t_m) \in \Theta}$.
\end{enumerate}
$\Struct{M}$ iku \emph{modhèl tèrma} amarga domainé
himpunan tèrma, lan !!{constant}s sarta !!{function}s
ditafsiraké saéngga $\Value{t}{M} = t$. !!{formula}s
atomik ing~$\Theta$ digunakaké kanggo netepaké
$\Assign{R}{M}$ saéngga $\Sat{M}{R(t_1, \dots, t_n)}$
yèn lan mung yèn $R(t_1, \dots, t_n) \in \Theta$.

\begin{lem}\ollabel{lem:truth}
Kanggo saben $!A \in \Theta \cup \Xi$:
\begin{enumerate}
  \item Yèn $!A \in \Theta$, mula $\Sat{M}{!A}$.
  \item Yèn $!A \in \Xi$, mula $\Sat/{M}{!A}$.
\end{enumerate}
\end{lem}

\begin{proof}
  Kanthi induksi miturut $\depth{!A}$.

Sepisan, anggep $!A$ atomik, yaiku $!A \ident \lfalse$
utawa $!A \ident R(t_1, \dots, t_m)$.

Amarga $\Pi_n \Sequent \Lambda_n$ ana ing cabang gagal,
ora ana $\Pi_n$ kang ngemot $\lfalse$ (yèn ana,
$\Pi_n \Sequent \Lambda_n$ dadi aksioma).
$\Sat{M}{R(t_1, \dots, t_m)}$
yèn lan mung yèn $R(t_1, \dots, t_m) \in \Theta$
miturut définisi~$\Struct{M}$.
Saka loro prakara iki, yèn $!A \in \Theta$,
mula $\Sat{M}{!A}$.

Yèn $!A$ atomik lan $!A \in \Pi_n$, mula
$!A \in \Pi_{n+1}$; yèn $!A \in \Lambda_n$, mula
$!A \in \Lambda_{n+1}$. (Iki manut cara algoritma
panelusuran bukti ngasilaké sekuèn panerus.) Mula,
yèn $!A \in \Theta \cap \Xi$, ana~$n$ saéngga
$!A \in \Pi_n \cap \Lambda_n$. Iki ndadèkaké
$\Pi_n \Sequent \Lambda_n$ aksioma, padahal cabang gagal
ora bisa ngemot aksioma. Dadi,
$!A \notin \Theta \cap \Xi$ kanggo~$!A$ atomik
miturut konstruksi. Akibaté,
yèn $!A \in \Xi$, mula $!A \notin \Theta$.
Mula, yèn $!A \in \Xi$, iki ateges $\Sat/{M}{!A}$
miturut définisi~$\Struct{M}$.

Kanggo langkah induksi, anggep (1) lan (2) lumaku tumrap
kabèh !!{formula}s kang !!{depth} luwih cilik tinimbang~$!A$.
Kita mbuktèkaké (1) lan (2) tumrap $!A$ kanthi mbedakaké
kasus-kasusé.

Gatèkna dhisik: yèn $!A^i$ katon ing
$\Pi_n \Sequent \Lambda_n$, $!A^i$ uga katon ing
$\Pi_{n+1} \Sequent \Lambda_{n+1}$ kajaba yèn $i$
indeks paling cilik ing antarané rumus nonatomik ing
$\Pi_n \Sequent \Lambda_n$. Amarga indeks paling cilik
ing antarané rumus nonatomik ing sekuèn paling dhuwur
kang ditambahaké ing saben tataran saya mundhak, pungkasane
kita tekan $\Pi_n \Sequent \Lambda_n$ kang ngemot $!A^i$
lan $i$ dadi indeks rumus nonatomik paling cilik.
Ing tataran~$n+1$,
algoritma ngreduksi~$!A^i$. Tegesé, yèn $!A \in \Theta$,
ana $n$ saéngga $\Pi_n = \Pi_n', !A^i$ lan $i$
indeks rumus nonatomik paling cilik. Semono uga,
yèn $!A \in \Xi$,
ana $n$ saéngga $\Lambda_n = \Lambda_n', !A^i$
lan $i$ indeks rumus nonatomik paling cilik.

\begin{enumerate}
  \item $!A \in \Theta$ lan $!A \ident !B \land !C$.
  Mula kanggo sawijiné $n$,
  $\Pi_n = \Pi_n', (!B \land !C)^i$.
  $\Pi_{n+1} \Sequent \Lambda_{n+1}$ iku premis
  kang cocog saka~\LeftR{\land}, yaiku
  $\Pi_{n+1} = \Pi_n', !B^{k+1}, !C^{k+2}$.
  Akibaté, $!B, !C \in \Theta$. Miturut hipotèsis
  induksi, $\Sat{M}{!B}$ lan $\Sat{M}{!C}$.
  Saka définisi $\Sat{M}{}$, kita olèh
  $\Sat{M}{!B \land !C}$.

  \item $!A \in \Xi$ lan $!A \ident !B \land !C$.
  Mula kanggo sawijiné $n$,
  $\Lambda_n = \Lambda_n', !B \land !C$.
  $\Pi_{n+1} \Sequent \Lambda_{n+1}$ iku salah siji
  saka rong premis~\RightR{\land} kang konklusiné
  $\Pi_n \Sequent \Lambda'_n, !B \land !C$, yaiku
  $\Lambda_{n+1} = \Lambda_n', !B^{k+1}$ utawa
  $\Lambda_{n+1} = \Lambda_n', !C^{k+1}$.
  Akibaté, $!B \in \Xi$ utawa $!C \in \Xi$.
  Miturut hipotèsis induksi, $\Sat/{M}{!B}$ utawa
  $\Sat/{M}{!C}$. Saka définisi $\Sat{M}{}$,
  kita olèh $\Sat/{M}{!B \land !C}$.
  
  \item Kasus $!A \ident \lnot !B$,
  $!A \ident !B \lor !C$, lan
  $!A \ident !B \lif !C$ padha carané lan
  dipasrahaké dadi latihan.

  \item $!A \in \Theta$ lan
  $!A \ident \lexists[x][!B(x)]$. Mula kanggo sawijiné
  $n$, $\Pi_n = \Pi_n', \lexists[x][!B(x)]^i$.
  $\Pi_{n+1} \Sequent \Lambda_{n+1}$ iku premis
  kang cocog saka~\LeftR{\lexists}, yaiku
  $\Pi_{n+1} = \Pi_n', !B(c)^{k+1}$ kanggo sawijiné~$c$.
  Akibaté, $!B(c) \in \Theta$ lan~$c \in C$.
  Miturut hipotèsis induksi, $\Sat{M}{!B(c)}$.
  Saka définisi $\Sat{M}{}$, kita olèh
  $\Sat{M}{\lexists[x][!B(x)]}$.

  \item $!A \in \Xi$ lan
  $!A \ident \lexists[x][!B(x)]$. Mula kanggo sawijiné
  $n$, $\Lambda_n = \Lambda_n', \lexists[x][!B(x)]^i$.
  Saben $\lexists[x][!B(x)]$ direduksi,
  $\lexists[x][!B(x)]$ tetep ana ing sisih tengen sekuèn
  premis kanthi indeks anyar. Mula, yèn formula iku ana
  ing sisih tengen cabang gagal,
  $\lexists[x][!B(x)]$ direduksi tanpa wates ing sisih
  tengen. Saben tèrma $t \in \Domain{M}$ pungkasane dadi
  ``tèrma kapisan kang mung ngemot konstanta ing~$C_n$
  lan durung digunakaké ing reduksi tengen
  $\lexists[x][!B(x)]$'' ing cabang gagal. Mula,
  kanggo saben $t \in \Domain{M}$,
  $!B(t) \in \Lambda_n$ kanggo sawijiné~$n$,
  lan akibaté $!B(t) \in \Xi$.
  Miturut hipotèsis induksi, $\Sat/{M}{!B(t)}$
  kanggo saben $t \in \Domain{M}$. Saka définisi
  $\Sat{M}{}$, kita olèh
  $\Sat/{M}{\lexists[x][!B(x)]}$.

  \item Kasus $!A \ident \lforall[x][!B(x)]$
  padha carané lan dipasrahaké dadi latihan.
\end{enumerate}
\end{proof}

\begin{cor}\ollabel{cor:complete} Algoritma panelusuran bukti
kanggo \Log{G3c} lengkap: yèn mandheg tanpa asil utawa
ora tau mandheg, sekuèn wiwitan
$\Gamma \Sequent \Delta$ ora valid, mula ora nduwèni
!!{proof}.
\end{cor}

\begin{proof}
  Yèn algoritma mandheg tanpa asil utawa ora tau mandheg,
  ana cabang gagal kang bisa dadi dhasar netepaké
  $\Struct{M}$. Miturut \olref{lem:truth}, kanggo
  saben $!A \in \Gamma$, $\Sat{M}{!A}$, lan kanggo
  saben $!A \in \Delta$, $\Sat/{M}{!A}$.
  (Elinga yèn $\Pi_0 = \Gamma$ lan $\Lambda_0 =
  \Delta$, mula $\Gamma \subseteq \Theta$ lan
  $\Delta \subseteq \Xi$.) Dadi
  $\Gamma \Sequent \Delta$ ora valid. Amarga
  \Log{G3c} sahih, $\Gamma \Sequent \Delta$
  ora nduwèni !!{proof}.
\end{proof}

\begin{cor}
\Log{G3c} lengkap: yèn $\Gamma \Sequent \Delta$ valid,
$\Log{G3c} \Proves \Gamma \Sequent \Delta$.
\end{cor}

\begin{proof}
  Kita mbuktèkaké kontraposisiné. Anggep
  $\Log{G3c} \Proves/ \Gamma \Sequent \Delta$.
  Mula algoritma panelusuran bukti mesthi mandheg tanpa
  asil utawa ora tau mandheg, amarga ora ana !!{proof}
  kang bisa ditemokaké. Miturut \olref{cor:complete},
  $\Gamma \Sequent \Delta$ ora valid.
\end{proof}


\end{document}
